\chapter{Features}
\label{chap:features}

% Forest styles for architecture snippets used in this chapter
\forestset{
  default preamble={
    for tree={
    draw,
    rounded corners,
    node options={align=left, inner xsep=5pt, inner ysep=4pt},
    s sep=5mm,
      l sep=7mm,
      anchor=north,
      edge={-stealth, line width=0.6pt},
      fill=white
    }
  },
  smallnode/.style={
    text width=3.0cm,
    font=\small
  },
  compact/.style={
    text width=2.6cm,
    font=\footnotesize
  }
}
\section{Routing}
\subsection{Introduction}
The Routing Module constitutes the foundational architecture of the "Transportation Platform." Its primary objective is to resolve a critical inefficiency within Ho Chi Minh City's transit ecosystem: the current lack of a unified location search and multi-modal route calculation capability. 
This module is designed to process and deliver optimized, dynamic routes across all relevant transit modalities, including private vehicles, bus, and pedestrian options.

To ensure adherence to the accelerated 11-week development schedule, and, critically, to leverage the robust real-time traffic data necessary for meaningful route optimization, the module will be developed via the integration of TomTom's Search and Routing APIs. 

Architecturally, this module will operate as a decoupled backend service (internal API), exposing three primary functional endpoints for consumption by other platform components and team members.
\begin{itemize}
  \item \textbf{Geocoding/Search:} Convert text addresses into coordinates and vice versa.
  \item \textbf{Routing:} Calculate the optimal route (time, distance, list of coordinates) between two points.
  \item \textbf{Map Generation:} Automatically use data from the Routing step to create and export an interactive .html map file (using Folium), with the route, start point, and end point already drawn.
\end{itemize}

\subsection{Input \& Output}
Inputs: origin, destination, time, preferences. 

Outputs: candidate routes, costs, and explanations.

\subsection{Work Breakdown}      
Tasks are divided into 3 main functional groups, corresponding to the user's activity flow:
\subsubsection{Selecting Start and End Points (Search \& Geocoding)}
This is the step to get coordinates (Lat/Lon) from user input.
\begin{itemize}
    \item \textbf{Build Endpoint POST /api/search:}
    \begin{itemize}
        \item Receive address (text string) from the client.
        \item Integrate search\_api logic into the endpoint.
    \end{itemize}
    \item \textbf{Call TomTom Search API:}
    \begin{itemize}
        \item Handle general search logic (/search/{address}.json).
        \item Handle nearby search logic (/nearbySearch/.json) if the client provides the current location's lat, lon.
    \end{itemize}
    \item \textbf{Process results:}
    \begin{itemize}
        \item Filter results not belonging to "Ho Chi Minh".
        \item Standardize returned JSON: Only send necessary information to the UI (name, address, lat, lon).
    \end{itemize}
    \item \textbf{Security:} Hide the TomTom API key (do not expose to the frontend).
\end{itemize}

\subsubsection{Route Calculation (Routing)}
This is the step to calculate the path after having 2 sets of coordinates.
\begin{itemize}
  \item \textbf{Build Endpoint POST /api/route:}
    \begin{itemize}
        \item Receive \texttt{start\_lat}, \texttt{start\_lon}, \texttt{end\_lat}, \texttt{end\_lon} from the client.
        \item Receive options: \texttt{travelMode} (car, motorcycle...) and \texttt{routeType} (fastest, shortest...).
    \end{itemize}
  \item \textbf{Integrate \texttt{tomtom\_route} logic:}
    \begin{itemize}
        \item Call TomTom \texttt{calculateRoute} API.
        \item Pass the correct parameters (vehicle, route type, \texttt{traffic=true}).
    \end{itemize}
  \item \textbf{Error Handling:}
    \begin{itemize}
        \item Manage cases where TomTom doesn't find a route and returns 4xx, 5xx errors.
        \item Return a 404 error code to the UI if no route is found (\texttt{''routes'' not in data}).
    \end{itemize}
\end{itemize}

\subsubsection{Map Generation and Export}
This is the final step and the main output of the Routing Module.

After collecting and processing all data from the TomTom API, the module will be responsible for aggregating and drawing this data into a standalone interactive .html map file.

This task uses logic from the visualize\_path\_on\_map function in the demo code.

\begin{itemize}
  \item \textbf{Aggregate input data:}
  \begin{itemize}
    \item Receive \texttt{coords} (list of route coordinates).
    \item Receive \texttt{start} (start point coordinates).
    \item Receive \texttt{end} (end point coordinates).
  \end{itemize}
  \item \textbf{Render the map (Using Folium):}
  \begin{itemize}
    \item Initialize a \texttt{folium.Map} object and automatically fit the frame (\texttt{m.fit\_bounds}) to the route.
    \item Draw Markers: Add 2 Markers (''Start'' and ''End'' icons) at the \texttt{start} and \texttt{end} points.
    \item Draw the path (PolyLine): Draw the route on the map (including the border and the main line as in the demo code).
  \end{itemize}
  \item \textbf{Export file (Main task):}
  \begin{itemize}
    \item Call the \texttt{m.save()} function to save the entire map object (markers, path) as a single HTML file (e.g., \texttt{route\_map.html}) on the server's disk.
  \end{itemize}
  \item \textbf{Completion:}
  \begin{itemize}
    \item When the \texttt{route\_map.html} file has been successfully created and saved, the Routing Module's task is considered complete.
    \item The API (Endpoint \texttt{/api/route}) will return a simple confirmation message, e.g., {''\texttt{status}'': ''\texttt{success}'', ''\texttt{message}'': ''\texttt{Map created}: \texttt{route\_map.html}''}.
  \end{itemize}
\end{itemize}



\subsection{Requirements}
\subsubsection{Hardware}
\begin{itemize}
    \item Development Machine: 1 Personal Laptop (Windows, macOS, or Linux) with an internet connection for programming and API testing.
    \item Deployment Server: Cloud hosting service (such as Render) that supports Python backend applications.
\end{itemize}
This platform will host the FastAPI/Flask internal API and provide a public endpoint for integration with other modules.
\subsubsection{Software}
\begin{itemize}
    \item \textbf{Language \& Framework:}
    \begin{itemize}
        \item Language: Python 3.x.
        \item API Framework: FastAPI (recommended for high performance and automatic API documentation) or Flask (simple, easy to set up).
    \end{itemize}
    \item \textbf{Main Python Libraries (Dependencies):}
    \begin{itemize}
        \item requests: Essential library for sending HTTP requests (API calls) to the TomTom server.
        \item folium: Key library to render GPS data (coordinates, paths) and create interactive .html map files.
        \item urllib: Use the urllib.parse submodule to encode URLs before calling the API.
        \item fastapi: Framework library for building the API.
        \item uvicorn: ASGI server, necessary to run the FastAPI application.
        \item python-dotenv: Used to manage the TomTom API key securely in an .env file.
    \end{itemize}
    \item \textbf{Services \& APIs Used (External Services):}
    \begin{itemize}
        \item TomTom Developer Account: Mandatory for registering and receiving an API Key.
        \item TomTom Search API: API used for the search function, converting addresses (strings) to coordinates (Lat/Lon).
        \item TomTom Nearby Search API: Used to search for locations within a 5km radius.
        \item TomTom Routing API: Core API, used to calculate routes, returns a summary (time, distance) and points (list of coordinates) to draw on the map.
    \end{itemize}
    \item \textbf{Development Tools \& Environment:}
    \begin{itemize}
        \item IDE/Editor: VS Code (Visual Studio Code).
        \item Version Control: Git and GitHub (to store code and collaborate).
        \item API Testing: Postman (an extremely important tool for testing the /api/search and /api/route endpoints I build before integration).
        \item Environment: Python Virtual Environment (venv) to manage project libraries in isolation.
    \end{itemize}
\end{itemize}


\subsection{Risk Management / Testing Method}

\begin{table}[H]
  \centering
  \renewcommand{\arraystretch}{1.15}
  \begin{tabularx}{\textwidth}{|>{\raggedright\arraybackslash}p{3.2cm}|>{\raggedright\arraybackslash}X|>{\raggedright\arraybackslash}p{1.8cm}|>{\raggedright\arraybackslash}X|}
    \hline
                   					\textbf{Risk} & \textbf{Description} & \textbf{Level} & \textbf{Mitigation Plan} \\
    \hline
    Cost / API Limits & TomTom provides 2{,}500 free calls/day. If usage exceeds this limit, the service will be suspended. & High & 1. Closely monitor the number of daily API calls.\newline 2. Research alternative plans such as using OpenStreetMap data. \\
    \hline
    Third-party Dependency & Service errors, maintenance, or slow response times from TomTom will paralyze the project's main functionality. & Medium & 1. Implement a reasonable timeout/retry mechanism.\newline 2. Clear error handling to notify the user. \\
    \hline
    Accuracy & Data in HCMC may not be 100\% accurate (small alleys, one-way streets, new roads). & Medium & 1. Accept as a provider limitation.\newline 2. Apply a filter to only keep results within HCMC, increasing relative accuracy. \\
    \hline
    Feature Limitations & Limited by the routeTypes TomTom provides. Cannot customize deeply (e.g., ``fewer red lights'', ``fewer lane changes''). & Low & 1. Adjust the project scope and accept the standard route types (fastest, shortest, balanced, eco). \\
    \hline
  \end{tabularx}
  \caption{Routing risks and mitigations}\label{tab:routing-risks}
\end{table}

\subsection{Timeline}
\renewcommand{\arraystretch}{1.15}
\begin{longtable}{|>{\raggedright\arraybackslash}p{0.12\textwidth}|>{\raggedright\arraybackslash}p{0.18\textwidth}|>{\raggedright\arraybackslash}p{0.34\textwidth}|>{\raggedright\arraybackslash}p{0.30\textwidth}|}
\caption{Routing timeline and milestones}\label{tab:routing-timeline}\\
\hline
    \textbf{Week} & \textbf{Phase} & \textbf{Main Task (Routing Module)} & \textbf{Objective / Checkpoint} \\
\hline
\endfirsthead
\hline
    \textbf{Week} & \textbf{Phase} & \textbf{Main Task (Routing Module)} & \textbf{Objective / Checkpoint} \\
\hline
\endhead
\hline
\multicolumn{4}{r}{Continued on next page} \\
\hline
\endfoot
\hline
\endlastfoot
Week 1 & A. Design & Requirements Analysis \& Scope: \newline\quad -- Finalize module scope. \newline\quad -- Team meeting to define integration points. & Check: Team agrees on the implementation scope as required. \\
\hline
Week 2 & A. Design & Technical Research: \newline\quad -- Compare API services (TomTom, Google\ldots). \newline\quad -- Begin researching the Folium library. & Check: Decision made to use TomTom API. \\
\hline
Week 3 & A. Design & API Service Design: \newline\quad -- Finalize 3 functions (Search, Route, Map). \newline\quad -- Share API design (Input/Output) with the team. & Check: Python test script for 3 functions: search, route, map. \\
\hline
Week 4 & B. Dev (POC) & Build Proof of Concept (POC): \newline\quad -- Start writing Python script. \newline\quad -- Integrate TomTom Search \& Route. & Check: Script runs and fetches JSON data from TomTom. \\
\hline
Week 5 & B. Dev (POC) & Complete POC: \newline\quad -- Integrate Folium (WBS 5.3). \newline\quad -- Filter ``Ho Chi Minh'' results. & Check: Code file (.py) works, successfully exports route\_map.html. \\
\hline
Week 6 & B. Dev & ``API-fication'' (Part 1): \newline\quad -- Install FastAPI, Uvicorn, python-dotenv. \newline\quad -- Move search\_api logic (WBS 5.1) into /api/search. & Check: Endpoint /api/search (on local) works, returns standard JSON. \\
\hline
Week 7 & B. Dev & ``API-fication'' (Part 2): \newline\quad -- Move tomtom\_route (WBS 5.2) and visualize\_path\_on\_map (WBS 5.3) logic into /api/route. & Check: Endpoint /api/route (on local) works, returns HTML file. \\
\hline
Week 8 & C. Integration & Deploy \& Integration: \newline\quad -- Deploy the FastAPI/Flask API to a cloud service that supports Python backend, such as Render. \newline\quad -- Start supporting the UI/UX Module to embed the interactive map using \texttt{<iframe>} with the public API URL from Render. & Check: Team has a public API URL. UI can display the map (even with bugs). \\
\hline
Week 9 & C. Integration & Error Handling \& Integration: \newline\quad -- Handle errors (Error Handling 4xx, 5xx). \newline\quad -- Support Safety Module \& Bus Route Finding System in fetching data. & Check: All other modules can call the API successfully. \\
\hline
Week 10 & D. Deploy & Testing \& Optimization: \newline\quad -- Fix final integration bugs. \newline\quad -- Consider caching (if needed). Finalize documentation. & Check: Functionality is stable in the test environment. \\
\hline
Week 11 & D. Deploy & Finalization \& Demo: \newline\quad -- Ensure API is stable on Production. \newline\quad -- Prepare slides \& demo. & Check: Ready for the final presentation. \\
\hline
\end{longtable}

\section{Recommendation System}
\subsection{Introduction}
Public transportation in Ho Chi Minh City is widely available, yet difficult for first-time visitors to navigate efficiently.

While mobile applications such as Google Maps and BusMap provide bus route guidance, they often lack transparency and flexibility in algorithmic design.

This project aims to build an independent system that helps visitors choose the optimal bus route from a starting stop to a destination stop, considering travel time, transfers, and waiting penalties.

However, a key challenge is the lack of open-source and up-to-date bus route data.

To address this, data was collected manually from official city transport sources and combined into structured CSV files, forming the foundational dataset of this project.
\subsection{Input \& Output}
\subsubsection{Input}
\begin{itemize}
    \item User Input:
    \begin{itemize}
        \item Start coordinate (latitude, longitude)
        \item Destination coordinate (latitude, longitude)
        \item MaxWalkingDistance (meter) {default = 300 m}: the maximum distance the user is willing to walk.
        \item MaxFares (VND) {default = infinity}: the maximum bus fares the user is willing to pay.
        \item Routing\_mode: cost-efficient / time-efficient
    \end{itemize}
    \item Dataset Files:
    \begin{itemize}
        \item \texttt{stops.csv} - list of bus stops with coordinates.
        \item \texttt{routes.csv} - metadata for each bus route (ID, name, fare, active time, spacing penalty).
        \item \texttt{trips.csv} - ordered sequence of stops for each route and distance between them.
        \item \texttt{car\_data.json} - used for visualizing bus paths on the map.
        \item \texttt{walk\_data.json} - used for visualizing walking paths on the map.
    \end{itemize}
\end{itemize}

\subsubsection{Output}
\begin{itemize}
    \item The optimal route (sequence of bus stops and route IDs) from start to destination.
    \item The estimated travel distance, estimated best scenario time, estimated worst scenario time (including waiting and transfer penalties).
    \item The total fares of the trip.
    \item  Visualized route on a map (differentiated each bus transfer).
\end{itemize}


\subsection{Work Breakdown}

\subsubsection{Data Layer}


\paragraph{Graph Construction}
\begin{itemize}
  \item Each bus stop is treated as a \textbf{node}.
  \item Adjacent stops on the same route are connected by \textbf{directed edges}.
  \item Edge weights represent \textbf{travel distance} (computed with the Haversine formula).
  \item A \textbf{transfer penalty} is added when switching between routes.
  %\item A \textbf{starting spacing penalty} is applied at the beginning to simulate real-world waiting time for the first bus.
\end{itemize}

\subsubsection{Pathfinding Layer}
Implements a modified $A^*$ search algorithm that considers:
\begin{itemize}
  \item Distance between stops.
  \item Transfer penalties between routes.
  \item Starting penalty for initial waiting time. \\
    \text{cost} = \text{currentcost} + \text{edgeDist}/$v$ + \text{transferPen} + c $\times$ \text{fare} + \text{startingPenalty}.

    \begin{itemize}
        \item currentcost -- accumulated cost to reach the current node.
        \item edgeDist -- distance (in meters) of the current bus segment.
        \item transferPen (s) -- transfer penalty, equal to the average bus spacing / waiting time; set transfer\_pen = 0 if there’s no transfer between routes.
        \item $v$ (m/s) -- average bus speed, assumed to be 7 m/s.
        \item $c$ -- constant coefficient that amplifies the importance of bus fare; set $c = 5$ if in cost-efficient mode, $c = 0$ if in time-efficient mode.
        \item fare (vnd) -- bus fare cost; fare = 0 if there’s no route transfer.
    \end{itemize}
    Estimated worst scenario time = TotalCost - TotalFares \\
    Estimated best scenario time = TotalCost - TotalFares - TotalTransferPenalties \\
    TotalDistance = estimated best scenario time $\times v$ + TotalWalkingDistance
\end{itemize}
The algorithm can compare multiple potential starting stops and select the route with the lowest total cost.

\subsubsection{Realistic Modeling Features}
\begin{itemize}
  \item \textbf{Real-time input:} Automatically remove bus routes not on active time.
  \item \textbf{Transfer Penalty:} Adds time when changing from one route to another, reflecting bus waiting time.
  \item \textbf{Starting Penalty:} Models the time a visitor might wait for the first bus to arrive.
  \item \textbf{Heuristic Function:} Uses geographical distance (via the Haversine formula) to guide the $A^*$ search efficiently.
\end{itemize}


\subsection{Requirements}
\subsubsection{Hardware}
\begin{itemize}
    \item Development Machine: 1 Personal Laptop (Windows, macOS, or Linux) with an internet connection for programming and API testing.
    \item Deployment Server: Cloud hosting service (such as Render) that supports Python backend applications.
\end{itemize}
This platform will host the FastAPI/Flask internal API and provide a public endpoint for integration with other modules.

\subsubsection{Software}
\begin{itemize}
  \item \textbf{Language:} Python (3.10+ recommended) as the primary implementation language.
  \item \textbf{Data Handling:} \texttt{pandas} and Python's \texttt{csv} module for reading, validating, and transforming bus stop and route data from CSV files (e.g., \texttt{stops.csv}, \texttt{routes.csv}, \texttt{trips.csv}).
  \item \textbf{Graph Algorithms:} Custom implementations of A* and Dijkstra with a pluggable cost function that supports penalties for waiting time, number of transfers, and walking distance; must return the path (sequence of stop IDs and route IDs) and aggregate cost/time metrics.
  \item \textbf{Math \& Geo Utilities:} Standard \texttt{math} package and a Haversine formula utility for great-circle distance between coordinates; define a small \emph{epsilon} tolerance for coordinate comparisons.
  \item \textbf{Visualization:} \texttt{folium} to render routes and stops on an interactive web map and export to a standalone HTML file for demos/presentations.
  \item \textbf{Documentation:} A Markdown report (\texttt{report.md}) describing dataset schemas, algorithm design/parameters, CLI usage examples, and known limitations.
\end{itemize}


\subsection{Timeline}
\begin{table}[H]
  \centering
  \renewcommand{\arraystretch}{1.15}
  \begin{tabularx}{\textwidth}{|>{\raggedright\arraybackslash}p{3.2cm}|>{\raggedright\arraybackslash}p{3.2cm}|>{\raggedright\arraybackslash}X|}
    \hline
    	\textbf{Week} & \textbf{Task} & \textbf{Description} \\
    \hline
    	\textbf{Week 1 -- Week 2} & \textbf{Research and Problem Definition} & Study existing transport apps, define goals and system constraints. Identify the absence of open-source bus data. \\
    \hline
    	\textbf{Week 3 -- Week 4} & \textbf{Data Collection and Cleaning} & Manually collect bus stop, route, and trip data from official or public sources. Convert them into standardized CSV format (\texttt{stops.csv}, \texttt{routes.csv}, \texttt{trips.csv}). \\
    \hline
    	\textbf{Week 5 -- Week 6} & \textbf{Graph Construction} & Implement logic to parse CSV files, create nodes and weighted edges, and structure the graph for routing (directed edges between adjacent stops, Haversine-based weights, transfer/start penalties). \\
    \hline
      \textbf{Week 7 -- Week 8} & \textbf{Algorithm Development} & Implement Dijkstra and $A^*$ algorithms with transfer and spacing penalties. Validate correctness with sample cases and unit tests. \\
    \hline
    	\textbf{Week 9 -- Week 10} & \textbf{Testing and Optimization} & Test with multiple start/destination combinations. Fine-tune penalties and heuristic to ensure realistic results and acceptable performance. \\
    \hline
    	\textbf{Future Work} & \textbf{Integration and Visualization} & Build a simple UI or map visualization to display the optimal bus route and estimated travel time to users (folium/OSM integration, optional lightweight web frontend). \\
    \hline
  \end{tabularx}
  \caption{Recommendation system timeline and milestones}\label{tab:reco-timeline}
\end{table}

\section{Safety \& Monitoring}
\subsection{Introduction}
HCMC endless traffic congestions especially in rush hour cannot be solved easily as the city is growing rapidly in both population and vehicles. 

To help commuters avoid traffic jams, we propose a traffic-aware routing feature that leverages real-time traffic camera feeds to analyze current traffic conditions and suggest alternative route. 

By integrating computer vision techniques with graph-based routing algorithms, this feature aims to enhance urban mobility and reduce travel times for users navigating through Ho Chi Minh City.
\paragraph{Impact}
\begin{itemize}
  \item Inefficient planning and delays.
  \item Poor commuter satisfaction.
\end{itemize}   

\subsection{Input \& Output}
\subsubsection{Input}
\begin{itemize}
  \item Start location.
  \item End location.
\end{itemize}

\paragraph{Output}
\begin{itemize}
  \item Coloured route according to the current traffic flow.
  \item Suggest one alternative less crowded route (if not possible, display a notice).
  \item A cute note for users.
\end{itemize}

\subsection{Work Breakdown}
\paragraph{Data \& Telemetry Pipeline}

\mbox{}\newline
	\textbf{Ingestion}
\begin{itemize}
  \item Pull JPG snapshots every 30 minutes from official HCMC traffic camera endpoints (\href{https://giaothong.hochiminhcity.gov.vn/Map.aspx}{city traffic portal}).
  \item De-duplicate by perceptual/image hash to avoid storing identical frames.
  \item Bind metadata: \texttt{camera\_id}, \texttt{segment\_id}, GPS, bearing, lane schema, timestamp.
  \item Persist raw frames + minimal JSON into object storage (Bronze layer).
\end{itemize}
  \textbf{Curation Layers}
\begin{itemize}
  \item Bronze: raw frames + minimal JSON (timestamp, camera\_id).
  \item Silver: per-frame detections, tracks, lane associations, calibration records, COCO annotations (labeled subset).
  \item Gold: per-segment time series every 1 hour: density (veh/km/lane), occupancy, queue tail, space-mean speed, incident flags.
\end{itemize}
	\textbf{Labeling Strategy}
\begin{itemize}
  \item Seed with 20--30k annotated frames (COCO boxes; classes: four-wheel vehicles, motorbike).
  \item Active learning to select ``hard'' frames (night rain, glare, heavy occlusion); add 1--2k frames per corridor during pilot.
\end{itemize}
	\textbf{Per-camera Calibration \& Geo-projection}
\begin{itemize}
  \item Use road corners to compute homographies.
  \item Back-project pixel distances to meters for speed/queue length estimation (standard monocular stationary-camera procedure).
\end{itemize}

\paragraph{Perception Stack}
\mbox{}\newline
\textbf{Pre-processing}
\begin{itemize}
  \item Defog/dehaze and gamma correction as first-stage normalization (stabilize detectors across dusk/night and haze).
  \item Legacy self-adaptive background subtraction with dynamic thresholds as secondary fallback for heavy occlusion/static cams.
\end{itemize}
	\textbf{Detection}
\begin{itemize}
  \item YOLO family (anchor-free head, e.g., \texttt{YOLOv11-S}); train with augmentation (mosaic, copy-paste), focal loss, finer grids for motorbikes.
\end{itemize}
	\textbf{Tracking \& Lane Association}
\begin{itemize}
  \item ByteTrack or OC-SORT at 2--5 FPS for ID-consistent tracks.
  \item Per-track lane assignment via calibrated lane polygons.
  \item Compute space-mean speed from track displacement in world coordinates.
\end{itemize}
	\textbf{Segment-level Signals (aggregation window $\Delta t = 2\text{--}4\ \text{min}$)}
\begin{itemize}
  \item Density (veh/km/lane), occupancy, speed percentiles, stop-and-go wave indicators; aggregate into Gold time series.
\end{itemize}

\paragraph{Short-horizon Traffic Prediction (5--30 min)} 


\textbf{Problem Framing}
\begin{itemize}
  \item Predict per-segment congestion class (R/Y/G) and/or travel time / mean speed at 5, 10, 20, 30 min ahead.
  \item Inputs: recent Gold signals + neighbor segment features.
\end{itemize}
	\textbf{Models}
\begin{itemize}
  \item Champion spatio-temporal GNN: DCRNN / STGCN / GraphWaveNet over road graph (nodes = segments, edges = adjacency/turns).
\end{itemize}
	\textbf{Feature Set \& Horizons}
\begin{itemize}
  \item Sliding windows: last 30--90 min at 2--4 min resolution; lags, moving averages, trend features.
  \item Exogenous: hour/day/week, holiday indicator, weather (rain intensity, temperature), seasonality.
\end{itemize}
	\textbf{Losses \& Metrics}
\begin{itemize}
  \item Regression (speed/ETA): RMSE, MAE, $R^2$.
  \item Classification (R/Y/G): macro-F1, precision, recall; horizon-specific evaluation ($\le 30$ min is primary use-case).
\end{itemize}

\paragraph{Time-dependent Routing}
\mbox{}\newline

	\textbf{Graph \& Costs}
\begin{itemize}
  \item Build time-dependent directed graph from OSM + city topology.
  \item Edge cost = predicted travel time $\times$ congestion penalty multiplier (e.g., Red 1.35, Yellow 1.15) + uncertainty surcharge when predictor confidence is low.
\end{itemize}
	\textbf{Search}
\begin{itemize}
  \item A* with landmark heuristics; do not exceed shortest-path ETA by >10\% unless red-jam risk predicted in next 15 min; avoid flooded links under weather alerts.
\end{itemize}
	\textbf{UX}
\begin{itemize}
  \item Show route colored by predicted flow.
  \item Provide one named alternative with rationale: ``Queue forming at Nguyễn Hữu Cảnh in 10 min; taking Điện Biên Phủ saves 7--10 min.''
  \item Cute note (copy layer): ``Sky is moody, roads are not---swing right to avoid the jam ahead.''
\end{itemize}

\paragraph{Privacy \& Safety}


\begin{itemize}
  \item No PII retention; mask plates/faces if visible.
  \item Edge redaction before storage; compliance with local data policies.
\end{itemize}


\paragraph{Ambiguities \& Limitations}
\begin{itemize}
  \item Motorbike occlusions in dense rain: under-count risk; mitigate via temporal smoothing + uncertainty flags (Medium limitation).
  \item Calibration drift (camera bumps; road works): re-fit homographies monthly or on drift alerts; residual errors affect speed more than counts (Low--Medium).
  \item Weather data latency/coverage: delayed granular rain data degrades prediction during storms; cache nearby station feeds and backfill (Low--Medium).
\end{itemize}

\subsection{Requirements}
\subsubsection{Hardware}
\begin{itemize}
  \item \textbf{Edge node (per camera or small group)}: CPU with 8--16 GB RAM; (optional GPU). Local SSD for short-term buffering (\textasciitilde24--48 h) and logs.
  \item \textbf{Central node (training/aggregation)}: Multi-core CPU, CUDA-capable GPU for fine-tuning detectors/GNNs; network-attached storage for Silver/Gold tables and artifacts.
  \item \textbf{Network}\,: Stable uplink from cameras or RTSP gateway; resilience via buffering and auto-retry on loss.
\end{itemize}

\subsubsection{Software}
\begin{itemize}
  \item \textbf{Perception}: YOLO family detector (\texttt{YOLOv8/YOLOv11-S}), ByteTrack/OC-SORT tracker, homography-based calibration utilities.
  \item \textbf{Data pipeline}: Services to write Bronze/Silver/Gold layers; time-window aggregation (2--4 min); REST/WebSocket publisher.
  \item \textbf{Prediction}: Spatio-temporal GNN framework (DCRNN/STGCN/GraphWaveNet) and a tree-based fallback (RF/GBT).
  \item \textbf{Dashboard}: Web UI for live map, history, and metrics; role-based access; configuration for junction onboarding.
\end{itemize}

\subsubsection{Functional Requirements (FR)}
\renewcommand{\arraystretch}{1.15}
\begin{tabularx}{\textwidth}{|>{\raggedright\arraybackslash}p{1.1cm}|>{\raggedright\arraybackslash}p{3.2cm}|>{\raggedright\arraybackslash}X|}
\hline
	\textbf{ID} & \textbf{Requirement} & \textbf{Description} \\
\hline
FR1 & Detector Setup & YOLOv11 model with vehicle classes; junction fine-tuning. \\
\hline
FR2 & Tracking \& Counting & Multi-object tracker + virtual lines per lane/movement; de-dup via track IDs. \\
\hline
FR3 & Class-wise Flow & Aggregate counts per 10--15 s (configurable) by class \& movement. \\
\hline
FR4 & Camera Calibration & Homography mapping pixels $\rightarrow$ lanes/movements. \\
\hline
FR5 & Graph Builder & Construct directed graph (junctions/edges) with rolling features. \\
\hline
FR6 & GNN Predictor & Short-term congestion per edge using recent flows \& exogenous signals. \\
\hline
FR7 & Alerting & Map scores to R/Y/G with thresholds + hysteresis; publish via REST/WebSocket. \\
\hline
FR8 & Dashboard & Live map with colored segments; history \& metrics view. \\
\hline
FR9 & Evaluation & mAP (detector), MAE/MAPE (counts), P/R/F1 (alerts); periodic reports. \\
\hline
\end{tabularx}

\subsubsection{Non-Functional Requirements (NFR)}
\renewcommand{\arraystretch}{1.15}
\begin{tabularx}{\textwidth}{|>{\raggedright\arraybackslash}p{1.1cm}|>{\raggedright\arraybackslash}p{3.2cm}|>{\raggedright\arraybackslash}X|}
\hline
	\textbf{ID} & \textbf{Requirement} & \textbf{Description} \\
\hline
NFR1 & Latency & Edge detection+count $\le 300$--$600$ ms/frame; end-to-end alert $\le 3$--$5$ s. \\
\hline
NFR2 & Accuracy & Counting MAE $\le 10\%$ on audited windows; Alert F1 $\ge 0.85$ (initial). \\
\hline
NFR3 & Robustness & Operate across day/night \& rain; track degradation. \\
\hline
NFR4 & Privacy & No faces/plates stored; only aggregates persist. \\
\hline
NFR5 & Reliability & 99\% pipeline uptime; auto-restart \& buffering on network loss. \\
\hline
NFR6 & Maintainability & Modular services; config-driven junction onboarding. \\
\hline
\end{tabularx}

\subsection{Risk Management / Testing Method}
\begin{itemize}
	\item \textbf{Perception (offline)} mAP@50 / @50--95 by class; IDF1/HOTA for tracking; speed MAE/MAPE vs radar/loop ground truth.
	\item \textbf{Prediction (offline)} Cross-validated MAE/RMSE (speed/ETA); macro-F1/precision/recall (R/Y/G); analyze by horizon (5/10/20/30 min) and weather bins.
	\item \textbf{Shadow Mode (online)} Compare predicted vs realized congestion/ETA for 4--6 weeks; calibrate alert thresholds to reach $\ge 80\%$ precision before public alerts.
\end{itemize}

\subsection{Timeline}
\begin{table}[H]
  \centering
  \renewcommand{\arraystretch}{1.15}
  \begin{tabularx}{\textwidth}{|>{\raggedright}p{1.6cm}|>{\raggedright}p{3.7cm}|>{\raggedright\arraybackslash}X|}
    \hline
    \textbf{Week} & \textbf{Task} & \textbf{Description / Exit criteria} \\
    \hline
    \textbf{1} & \textbf{Project Setup \& Governance} &
    Repo, CI, coding style; success metrics (mAP/IDF1, MAE/F1), data policy; camera list \& schema freeze. Deliverable: 1–2 pp scope \& risk register. \\
    \hline
    \textbf{2} & \textbf{Ingestion Prototype (Bronze)} &
    Pull JPG snapshots from HCMC endpoints every 30 min; perceptual hash de-dup; bind metadata (camera\_id, GPS, bearing). Deliverable: 24 h Bronze sample. \\
    \hline
    \textbf{3} & \textbf{Label Seed \& Calibration Kit} &
    Curate 20–30k COCO boxes (car/motorbike); active-learning sampler; homography tool (road corners). Deliverable: 5 calibrated cams + seed set. \\
    \hline
    \textbf{4} & \textbf{Detector Baseline} &
    Train/tune YOLOv8/YOLOv11-S with augmentation; night/rain sanity checks. Exit: mAP@50 baseline achieved; error notebook for hard scenes. \\
    \hline
    \textbf{5} & \textbf{Tracking \& Counting} &
    ByteTrack/OC-SORT at 2–5 FPS; virtual lines per lane/movement; de-dup via track IDs. Exit: MAE\,$\le$\,10–15\% on audited windows. \\
    \hline
    \textbf{6} & \textbf{Silver/Gold Pipeline \& Metrics} &
    Aggregate $\Delta t=2$–4 min to density/occupancy/speed; persist Silver/Gold; dashboard skeleton. Exit: 1 week of Gold time series. \\
    \hline
    \textbf{7} & \textbf{Road Graph \& Features} &
    Build directed graph (OSM + topology); engineer features/lags; create train/val splits; tree baseline (RF/GBT) for ETA/speed. Exit: offline MAE/RMSE baseline. \\
    \hline
    \textbf{8} & \textbf{Spatio-temporal GNN} &
    Train DCRNN/STGCN/GraphWaveNet for 5/10/20/30-min horizons; calibrate uncertainty; map to R/Y/G with hysteresis. Exit: macro-F1 target for $\le$30 min. \\
    \hline
    \textbf{9} & \textbf{Time-dependent Routing (A*)} &
    Edge cost = predicted travel time $\times$ congestion penalty + uncertainty surcharge; one named alternative route + rationale; flood avoidance hook. Exit: prototype UI coloring. \\
    \hline
    \textbf{10} & \textbf{Shadow Mode \& Tuning} &
    Live compare predicted vs realized congestion/ETA; threshold calibration to reach $\ge$80\% precision before alerts; privacy audit (plate/face masking). \\
    \hline
    \textbf{11} & \textbf{Hardening \& Demo} &
    Latency profiling (edge $\le$300–600 ms/frame; E2E alert $\le$3–5 s); reliability tests (auto-retry/buffering); docs, ops runbook, demo with “cute note”. \\
    \hline
  \end{tabularx}
  \caption{Safety \& Monitoring module: 11-week timeline and milestones}
  \label{tab:safety-monitoring-timeline}
\end{table}



\section{Camera Ingestion Pipeline}
\label{sec:camera-ingestion}
\subsection{Introduction}
			\textbf{Project purpose}\\
Collect, normalize, and serve legally compliant real-time traffic camera frames to support downstream modules (object detection, flow estimation, congestion scoring, safety alerts). The pipeline must deliver low-latency, robust, and auditable frame ingestion while protecting privacy and meeting retention policies.

\vspace{0.5em}
			\textbf{Scope}
\begin{itemize}
  \item Continuous ingestion from many heterogeneous camera sources (varying resolution, FPS, protocols).
  \item Per-camera preprocessing (ROI, normalization), calibration for geometric reasoning, short-term storage for debugging, and health/metrics dashboards.
  \item API access for downstream consumers and operators, plus compliance documentation and privacy protections.
\end{itemize}

		    \textbf{Key goals}
\begin{itemize}
  \item Reliable frame delivery (target freshness SLA).
  \item Privacy and legal compliance for each feed.
  \item Observability and easy onboarding of new cameras.
\end{itemize}

\subsection{Input \& Output}
\subsubsection{Inputs}
\begin{itemize}
  \item Camera feed endpoints (RTSP/HTTP/HLS/etc.) and credentials.
  \item Camera metadata: \texttt{camera\_id}, URL, GPS, orientation, vendor, resolution, nominal FPS, tilt/azimuth.
  \item Permission/terms-of-use / proof of right to ingest.
  \item Per-camera ROI masks and calibration images (for homography).
  \item Central configuration (sampling rate, retention policy, alert thresholds).
\end{itemize}

\subsubsection{Outputs}
\begin{itemize}
  \item Sampled frames (JPEG/PNG) and short clips retained for 24--72 hours in object storage.
  \item Per-frame metadata records (timestamp, camera\_id, seq\_no, ingest\_latency, checksum) in Postgres.
  \item Calibration artifacts (homography matrices, ROI masks) stored alongside camera metadata.
  \item Time-series metrics (frame freshness, fetch rate, errors, latency) in TimescaleDB/InfluxDB and Grafana dashboards.
  \item HTTP/REST API endpoints returning latest frames, frame metadata, and camera health.
  \item Audit \& compliance report (camera permissions, retention logs, redaction events).
\end{itemize}

\subsection{Work Breakdown}
\subsubsection{Camera Onboarding \& Registry}
\begin{itemize}
  \item Build camera registry schema (Postgres): metadata, legal status, sampling policy.
  \item CLI / UI to add/update cameras and upload ROI/calibration images.
\end{itemize}

\subsubsection{Scraper \& Fetcher}
\begin{itemize}
  \item Implement resilient fetcher supporting RTSP/HLS/HTTP with backoff, retries, rate limiting.
  \item Health checks, circuit breakers, and opt-in buffering for high-latency feeds.
\end{itemize}

\subsubsection{Frame Extraction \& Sampling}
\begin{itemize}
  \item FFmpeg integration to sample frames at configurable rates (default 1 fps) and normalize containers.
  \item Edge cases handling (variable FPS, partial frames).
\end{itemize}

\subsubsection{Preprocessing \& Calibration}
\begin{itemize}
  \item Image normalization (resize, color/contrast adjustment) using OpenCV.
  \item Apply ROI masks, blur redaction zones if required.
  \item Compute / store homography for selected cameras for downstream speed/trajectory estimation.
\end{itemize}

\subsubsection{Storage \& Retention}
\begin{itemize}
  \item Short-term frame storage: S3/MinIO with lifecycle rules (24--72 h).
  \item Persistent metadata in Postgres and time-series metrics in TimescaleDB/InfluxDB.
\end{itemize}

\subsubsection{API \& Access}
\begin{itemize}
  \item FastAPI endpoints: register/list cameras, latest frame, health, metrics.
  \item AuthN/AuthZ for operators and downstream services.
\end{itemize}

\subsubsection{Monitoring, Alerts \& Observability}
\begin{itemize}
  \item Export metrics (Prometheus) and dashboards (Grafana) for frame freshness, success rate, latency.
  \item Alert rules (stale frames, high error rate, permission expiry).
\end{itemize}

\subsubsection{Compliance \& Privacy}
\begin{itemize}
  \item Legal review workflow per camera; store permission evidence.
  \item Automatic face/license-plate blurring or ROI redaction when required.
  \item Retention enforcement and audit logs.
\end{itemize}

\subsection{Requirements}
\subsubsection{Functional Requirements}
\begin{enumerate}
  \item Register and manage camera metadata and legal status.
  \item Continuously fetch and sample frames from registered cameras.
  \item Provide per-camera ROI masking and optional redaction.
  \item Expose API to retrieve latest frames and camera health.
  \item Store frames for configurable short retention window and metadata indefinitely.
\end{enumerate}

\subsubsection{Non-Functional Requirements}
\begin{itemize}
  \item \textbf{Latency}: end-to-end frame availability within target (e.g., $<2\,\text{s}$ after sample time for on-network cameras).
  \item \textbf{Scalability}: support $N$ concurrent camera feeds (design target e.g., 1k) with horizontal scaling.
  \item \textbf{Reliability}: per-camera success rate target (e.g., $>99\%$ over rolling 24 h).
  \item \textbf{Security}: encrypted in-transit and at-rest; credentials stored safely.
  \item \textbf{Observability}: metrics for freshness, errors, and throughput; logs for forensic audits.
\end{itemize}

\subsubsection{Legal \& Compliance}
\begin{itemize}
  \item Track and store permission documents for each camera feed.
  \item Implement automatic redaction if terms require (faces, plates) and log redaction events.
  \item Enforce retention policy and provide exportable compliance evidence.
\end{itemize}

\subsubsection{Deployment \& Ops}
\begin{itemize}
  \item Containerized services (Docker) with CI/CD (GitHub Actions).
  \item Infrastructure as code (Terraform) for environments and S3/MinIO provisioning.
  \item Runbooks for onboarding and incident response.
\end{itemize}

\subsection{Risk Management / Testing Method}
\subsubsection{Major Risks \& Mitigations}
\begin{itemize}
  \item \textbf{Feed instability / network latency}: implement retry/backoff, local buffering, and fallback sampling frequency.
  \item \textbf{Heterogeneous stream formats / codecs}: standardize via FFmpeg transcoding layer and build format compatibility tests.
  \item \textbf{Privacy / legal exposure}: require documented permissions before enabling a camera; automatic redaction; legal review gate.
  \item \textbf{Scaling limits under high concurrency}: autoscaling policies, connection pooling, and rate limiting per-camera.
  \item \textbf{Poor image quality (weather, night)}: flag low-quality frames and route to quality-aware downstream models or human review.
\end{itemize}

\subsubsection{Testing Strategy}
\begin{enumerate}
  \item \textbf{Unit tests}: parser, registry CRUD, config validation.
  \item \textbf{Integration tests}: end-to-end sampling from a synthetic stream, storage round-trip, metadata writes.
  \item \textbf{Contract tests}: API schema/response shape for downstream consumers.
  \item \textbf{Load \& Stress tests}: simulate target concurrent feeds and max throughput (sampled frame rate $\times$ $N$ cameras).
  \item \textbf{Fault injection / Chaos testing}: simulate network drops, FFmpeg failure, permissions revocation to verify graceful recovery.
  \item \textbf{Security \& privacy tests}: verify redaction effectiveness (face/plate blur), permissions enforcement, token expiry.
  \item \textbf{Acceptance tests / E2E}: run a validation scenario covering registration $\rightarrow$ sampling $\rightarrow$ API retrieval $\rightarrow$ metric logging.
\end{enumerate}

\subsubsection{Acceptance Criteria}
\begin{itemize}
  \item Freshness SLA met for 95\%+ of feeds in baseline environment.
  \item Error rate below agreed threshold (e.g., $<1\%$ critical failures).
  \item Permissions recorded for every active camera; redaction applied where required.
  \item Monitoring dashboards and alerts firing for defined conditions.
\end{itemize}

\subsection{Timeline}
\renewcommand{\arraystretch}{1.15}
\begin{longtable}{|>{\raggedright\arraybackslash}p{0.14\textwidth}|>{\raggedright\arraybackslash}p{0.22\textwidth}|>{\raggedright\arraybackslash}p{0.56\textwidth}|}
\caption{Camera Ingestion Pipeline timeline and deliverables}\label{tab:camera-ingestion-timeline}\\
\hline
									\textbf{Week} & \textbf{Focus} & \textbf{Output / Deliverables} \\
\hline
\endfirsthead
\hline
									\textbf{Week} & \textbf{Focus} & \textbf{Output / Deliverables (cont.)} \\
\hline
\endhead
\hline
\multicolumn{3}{r}{\small Continued on next page} \\
\hline
\endfoot
\hline
\endlastfoot
Week 0 & Setup \& Legal & Camera inventory, permissions review, sampling policy \& test plan \\
\hline
Week 1--2 & Pipeline Core & Scraper, FFmpeg extractor, camera registry, basic health monitoring \\
\hline
Week 3--4 & Preprocessing & Frame normalization, ROI masking, redaction prototype \\
\hline
Week 5 & Calibration & Homography / perspective calibration for selected cameras \\
\hline
Week 6 & Metrics \& Logging & Frame freshness and uptime metrics stored in time-series DB \\
\hline
Week 7 & API Integration & Endpoints for latest frame, camera metadata, and health \\
\hline
Week 8 & Privacy Controls & Blurring/masking + retention enforcement implementation \\
\hline
Week 9 & Dashboards & Grafana monitoring dashboard and alert thresholds \\
\hline
Week 10 & E2E Testing & Full validation of pipeline stability, latency, and failure modes \\
\hline
Week 11 & Documentation & Runbooks, onboarding guide, compliance evidence packaging \\
\hline
\end{longtable}

\section{Database}
\label{sec:database}
\subsection{Introduction}
Persist, index, and serve all mobility data for routing, congestion prediction, and UX layers. Normalize heterogeneous sources (Overpass/OSM, GTFS-like CSVs, camera-derived signals) into time-aligned, queryable stores. Provide stable contracts (SQL/Parquet views + APIs) to upstream ML (GNN/TBMs) and downstream services (A* router, map UI).\\


\paragraph{Scope (what it does not own)}
\begin{itemize}
  \item No model training/inference logic (served by ML service).
  \item No image processing (YOLO/ByteTrack) --- Database only stores outputs \& metadata.
\end{itemize}


\subsection{Input \& Output}
\subsubsection{Inbound data (batch/stream)}
\begin{itemize}
  \item \textbf{Overpass/OSM snapshots (JSON)}: roads, nodes, ways; Python fetchers already in place.
  \item \textbf{City cameras} $\rightarrow$ per-frame detections/tracks $\rightarrow$ per-segment aggregates (density, queue length, speed proxies).
  \item \textbf{GTFS-like CSVs}: \texttt{routes.csv}, \texttt{stops.csv}, \texttt{trips.csv} (for stop proximity, segment naming, UX overlays).
  \item \textbf{Weather/seasonality features} (hour, weekday, rain flags) for ML enrichment.
\end{itemize}


\subsubsection{Canonical schemas (storage contracts)}
\paragraph{Bronze (immutable, as-landed)}
\begin{itemize}
  \item \texttt{raw\_overpass} (JSON blob + bbox, \texttt{fetched\_at}).
  \item \texttt{raw\_frames} (\texttt{uri}, \texttt{camera\_id}, \texttt{ts}, \texttt{img\_hash}).
  \item \texttt{raw\_detections} (\texttt{frame\_id}, \texttt{class}, \texttt{bbox}, \texttt{score}).
\end{itemize}

\paragraph{Silver (cleaned \& derived)}
\begin{itemize}
  \item \texttt{lane\_calibration} (\texttt{camera\_id}, \texttt{lane\_id}, \texttt{homography}, \texttt{updated\_at}).
  \item \texttt{tracks} (\texttt{track\_id}, \texttt{camera\_id}, \texttt{lane\_id}, \texttt{ts}, \texttt{world\_x}/\texttt{world\_y}, \texttt{speed\_est}, \texttt{is\_stalled}).
  \item \texttt{segments} (\texttt{segment\_id}, \texttt{geom}, \texttt{speed\_limit}, \texttt{lanes}, \texttt{osm\_way\_ids}).
\end{itemize}

\paragraph{Gold (time series for ML \& routing)}
\begin{itemize}
  \item \texttt{segment\_signals\_2m} (\texttt{ts}, \texttt{segment\_id}, \texttt{density\_veh\_per\_km\_lane}, \texttt{occupancy}, \texttt{speed\_p50}, \texttt{speed\_p85}, \texttt{queue\_tail\_m}, \texttt{incident\_flag}, \texttt{weather\_bins}).
  \item \texttt{congestion\_labels\_2m} (\texttt{ts}, \texttt{segment\_id}, \texttt{class} $\in \{\text{G},\text{Y},\text{R}\}$, \texttt{label\_method}, \texttt{confidence}).
  \item \texttt{travel\_time\_edge\_ETA\_2m} (\texttt{ts}, \texttt{edge\_id}, \texttt{eta\_s}, \texttt{eta\_conf}, \texttt{horizon\_min} $\in \{5,10,20,30\}$).
\end{itemize}

\subsubsection{Outbound interfaces}
	\textbf{SQL/Views for ML}: sliding-window tensors (RNN/GNN ready), normalized features (MinMax, one-hot).\\
	\textbf{HTTP/gRPC (internal)}:
\begin{itemize}
  \item \texttt{GET /db/segments/\{id\}/signals?from=\&to=}
  \item \texttt{GET /db/edges/eta?at=...\&horizon=...}
  \item \texttt{POST /db/congestion\_labels:write} (for online labels / human QA)
\end{itemize}

\subsection{Work Breakdown}
\subsubsection{Data Modeling \& Storage}
Define PostGIS schemas for segments, edges, stops, routes; Parquet for high-volume time series (signals, ETA). Partition by date (daily) + \texttt{segment\_id} bucketing; ZSTD compression.

\subsubsection{Ingestion Pipelines}
Overpass puller $\rightarrow$ \texttt{raw\_overpass}; Camera snapshot harvester $\rightarrow$ hash-dedupe $\rightarrow$ \texttt{raw\_frames}; detection/tracking appends to \texttt{raw\_detections}/\texttt{tracks}; CSV loaders for \texttt{routes.csv}, \texttt{stops.csv}, \texttt{trips.csv}.
\subsubsection{Transformation}
Build segments from OSM ways (respect oneway, service exclusions). Aggregate tracks $\rightarrow$ \texttt{segment\_signals\_2m} with calibrated lane polygons. Join exogenous features (weather, seasonality).

\subsubsection{Serving}
Materialized views for ML (sliding windows: last 30--90 min) and router (time-dependent edge costs). Low-latency KV (Redis/KeyDB) cache for “current” 0--30 min horizons.

\subsubsection{Governance \& Ops}
Data retention policy (Bronze 90 d; Silver/Gold 365 d rolling). CDC/versioning for calibration \& segment topology. Backfills + schema migration playbooks.

\subsection{Requirements}
\subsubsection{Functional}
\begin{itemize}
  \item Store $\ge 180$ days of \texttt{segment\_signals\_2m} for model retraining and seasonality.
  \item Provide $\le 300$ ms P95 reads for latest 30-minute window per corridor to serve A* routing.
  \item Offer labeled congestion classes (R/Y/G) + confidence for alerts (DT/RF fallback proven strong).
\end{itemize}

\subsubsection{Non-functional}
\begin{itemize}
  \item Reliability: 99.5\% monthly availability for reads; durable storage (daily backups).
  \item Latency: ingest $\rightarrow$ Gold availability $\le 90$ s after frame capture.
  \item Security/Privacy: no PII; plate/face hashes never stored.
\end{itemize}

\subsection{Risk Management / Testing Method}
\subsubsection{Key risks \& mitigations}
\begin{itemize}
  \item Camera downtime/data gaps $\rightarrow$ impute from neighbors + seasonality; flag low-confidence.
  \item Night/rain occlusion $\rightarrow$ temporal smoothing; fallback detectors.
  \item Calibration drift $\rightarrow$ version homographies, CDC rollback.
  \item OSM mismatch $\rightarrow$ maintain override layer + periodic diffs.
\end{itemize}

\subsubsection{Testing strategy}
\begin{itemize}
  \item Schema \& contract tests: migration dry-runs, backward-compatibility.
  \item Ingest fidelity tests: checksum \& row-count parity.
  \item Feature quality tests: drift monitors, missingness dashboards.
  \item Latency SLO tests: synthetic load (1k segments).
  \item Label QA: macro-F1 tracking (DT/RF baseline $\ge 0.90$).
\end{itemize}

\subsection{Timeline}
\renewcommand{\arraystretch}{1.15}
\begin{tabularx}{\textwidth}{|>{\raggedright\arraybackslash}p{1.2cm}|>{\raggedright\arraybackslash}p{3.2cm}|>{\raggedright\arraybackslash}X|>{\raggedright\arraybackslash}p{2.2cm}|}
\hline
	\textbf{Week} & \textbf{Focus} & \textbf{Deliverables} & \textbf{Certainty} \\
\hline
1--2 & Foundations & ERD, PostGIS/Parquet setup, Bronze/Silver/Gold scaffolding, CSV loaders & High \\
\hline
3--4 & Camera $\rightarrow$ Signals path & \texttt{raw\_frames}/\texttt{raw\_detections} landed, aggregation $\rightarrow$ \texttt{segment\_signals\_2m}, calibration versioning & Medium-High \\
\hline
5 & Gold \& Serving & Gold views, ML exports, \texttt{travel\_time\_edge\_ETA\_2m}, Redis cache & Medium \\
\hline
6 & QA \& Hardening & Quality monitors, backfills, retention/backup, load \& label QA loop & Medium \\
\hline
\end{tabularx}

\section{UI/UX Platform}
\label{sec:uiux}
\subsection{Introduction}
Urban travelers today face a fragmented mobility experience --- switching between multiple apps for bus schedules, ride-hailing, and walking navigation. This results in:
\begin{itemize}
  \item Inefficient trip planning (balancing time, cost, and distance).
  \item Lack of real-time actionable insights (traffic, ETA drift).
  \item User frustration and cognitive fatigue.
\end{itemize}

\paragraph{UI/UX Goal} Design a unified, intuitive, and responsive interface that allows users to:
\begin{itemize}
  \item Plan, book, and monitor their trips end-to-end in one platform.
  \item Receive real-time recommendations with clear trade-offs.
  \item Navigate safely with contextual alerts and visual clarity.
\end{itemize}

\subsection{Input \& Output}
\renewcommand{\arraystretch}{1.12}
\begin{tabularx}{\textwidth}{|>{\raggedright\arraybackslash}p{2.4cm}|>{\raggedright\arraybackslash}X|}
\hline
	extbf{Category} & \textbf{Description} \\
\hline
Input & Start \& destination points, preferred transport modes, live GPS data, transit APIs, local build API for object detection, and user preferences. \\
\hline
Output & Optimized travel itineraries (Fastest/Normal), ETA \& cost estimates, congestion alerts, and real-time route updates. \\
\hline
\end{tabularx}

\subsubsection{System Flow}
\begin{enumerate}
  \item User enters trip details.
  \item System fetches live transport \& traffic data \& camera data.
  \item UI displays top route options.
  \item In-trip recommender updates dynamically via GPS and APIs.
  \item Alerts and rerouting offered in case of congestion.
\end{enumerate}

\subsection{Work Breakdown}
\renewcommand{\arraystretch}{1.12}
\begin{tabularx}{\textwidth}{|>{\raggedright\arraybackslash}p{3.3cm}|>{\raggedright\arraybackslash}X|>{\raggedright\arraybackslash}p{3.6cm}|}
\hline
	extbf{Module} & \textbf{Description} & \textbf{Deliverables} \\
\hline
A. User Research \& Persona & Identify user pain points via surveys and interviews. & Persona profiles, user journey maps \\
\hline
B. Information Architecture & Design navigation across Pre-trip, In-trip, and Alerting modules. & Sitemap, navigation flow \\
\hline
C. Wireframe (Low-Fidelity) & Build layout skeletons for key screens (Trip Planner, Tracking). & Low-fi wireframes \\
\hline
D. Prototype (High-Fidelity) & Apply colors, typography, icons, and animations. & Interactive Figma prototype \\
\hline
E. UI Implementation & Convert designs to React/Next.js frontend with live API data. & Responsive working UI \\
\hline
F. UX Testing \& Iteration & Conduct usability tests and refine UI based on feedback. & UX test report, iteration summary \\
\hline
\end{tabularx}

\subsection{Requirements}
\subsubsection{Functional Requirements}
\begin{itemize}
  \item Allow users to input and edit trip details (start/end points, mode).
  \item Display optimized routes with ETA, cost, and traffic levels.
  \item Provide real-time updates and one-tap reroute options.
  \item Deliver congestion \& safety alerts.
  \item Ensure accessibility compliance (font size $\ge 14$px, contrast ratio $\ge 4.5:1$).
\end{itemize}

\subsubsection{Non-Functional Requirements}
\renewcommand{\arraystretch}{1.12}
\begin{tabularx}{\textwidth}{|>{\raggedright\arraybackslash}p{2.9cm}|>{\raggedright\arraybackslash}X|}
\hline
	\textbf{Aspect} & \textbf{Specification} \\
\hline
Usability & $\ge 85\%$ success rate in usability tests \\
\hline
Performance & Response latency $\le 5$ s (p95) \\
\hline
Reliability & $\ge 99.5\%$ crash-free sessions \\
\hline
Accessibility & WCAG 2.1 compliance \\
\hline
Consistency & Unified icons, color scheme, and feedback states \\
\hline
\end{tabularx}

\subsubsection{Tools \& Frameworks}
\renewcommand{\arraystretch}{1.12}
\begin{tabularx}{\textwidth}{|>{\raggedright\arraybackslash}p{3.0cm}|>{\raggedright\arraybackslash}X|}
\hline
	\textbf{Purpose} & \textbf{Tools / Frameworks} \\
\hline
UI Design & Figma, Adobe XD, Illustrator \\
\hline
Prototyping & Figma Prototype, Framer, InVision \\
\hline
Frontend & React.js, Next.js, Tailwind CSS, shadcn/ui \\
\hline
UX Testing & Maze, Hotjar, Google Analytics \\
\hline
Collaboration & Notion, Trello, Slack, GitHub \\
\hline
\end{tabularx}

\subsection{Risk Management / Testing Method}
\subsubsection{Risk Analysis}
\renewcommand{\arraystretch}{1.12}
\begin{tabularx}{\textwidth}{|>{\raggedright\arraybackslash}p{3.5cm}|>{\raggedright\arraybackslash}p{2.6cm}|>{\raggedright\arraybackslash}X|}
\hline
	\textbf{Risk} & \textbf{Impact} & \textbf{Mitigation Strategy} \\
\hline
API data inconsistency (ETA drift) & Medium & Cache fallback data; re-request on failure. \\
\hline
Overwhelming UI complexity & High & Simplify actions: $\le 3$ taps for core functions. \\
\hline
Accessibility issues & High & Use accessible color palette and test with screen readers. \\
\hline
Integration delays & Medium & Align weekly via shared Figma + GitHub sync. \\
\hline
Low participation in usability testing & Low & Utilize online testing via Maze/Hotjar. \\
\hline
\end{tabularx}

\subsubsection{Testing Methods}
\begin{itemize}
  \item \textbf{Usability Testing}: Evaluate ease of use and completion rate of key actions.
  \item \textbf{A/B Testing}: Compare alternative UI layouts for efficiency.
  \item \textbf{Heuristic Evaluation}: Check compliance with UX principles.
  \item \textbf{Performance Testing}: Measure app response time and smoothness.
  \item \textbf{User Satisfaction Survey}: Target CSAT $\ge 4.2/5$.
\end{itemize}

\subsection{Timeline}
\renewcommand{\arraystretch}{1.12}
\begin{tabularx}{\textwidth}{|>{\raggedright\arraybackslash}p{2.3cm}|>{\raggedright\arraybackslash}p{2.9cm}|>{\raggedright\arraybackslash}X|>{\raggedright\arraybackslash}p{3.4cm}|}
\hline
	\textbf{Week} & \textbf{Phase} & \textbf{Key Activities} & \textbf{Deliverables} \\
\hline
Week 1--2 & Research \& Persona & Conduct surveys/interviews and map user journeys. & Persona \& journey map \\
\hline
Week 3 & Information Architecture & Design sitemap and define module structure. & Sitemap \& flow diagram \\
\hline
Week 4--5 & Wireframing & Create low-fidelity wireframes for all key screens. & Approved wireframes \\
\hline
Week 6--7 & High-Fidelity Design & Apply brand system, icons, and UI interactions. & Interactive prototype \\
\hline
Week 8 & UI Implementation & Develop and integrate React/Next.js frontend. & Working UI with API link \\
\hline
Week 9--10 & UX Testing \& Iteration & Conduct usability tests, analyze feedback. & UX Report + Refined Design \\
\hline
Week 11 & Final Presentation & Prepare final demo and design deck. & Live demo \& presentation materials \\
\hline
\end{tabularx}
